\documentclass[11pt,a4paper]{article}
\usepackage[T1]{fontenc}
\usepackage[utf8]{inputenc}
\usepackage[french]{babel}
\frenchsetup{SmallCapsFigTabCaptions=false}
\usepackage{lmodern,microtype}
\usepackage[margin=23mm]{geometry}
\usepackage{booktabs,array,tabularx,enumitem,tikz}
\usepackage{hyperref}
\usetikzlibrary{arrows.meta,positioning}
\hypersetup{colorlinks=true,linkcolor=black,citecolor=black,urlcolor=black,
  pdftitle={LANCE : architecture du pipeline et méthodes de preuve},
  pdfauthor={Codex, agent de diagnostic statique},pdflang={fr-FR}}
\urlstyle{same}
\def\UrlBreaks{\do\/\do\-\do\_\do\.\do\:}
\newcolumntype{Y}{>{\raggedright\arraybackslash}X}
\newcommand{\code}[1]{\nolinkurl{#1}}
\newcommand{\src}[2]{{\small\href{../../#1}{\nolinkurl{#1}}, l.~#2}}
\newcommand{\revision}{2026-10-07 13:41 (Europe/Brussels, UTC+02:00)}
\setlength{\parskip}{3pt}
\setlength{\parindent}{1.2em}
\setlength{\emergencystretch}{2em}
\setlist{nosep,leftmargin=5mm,topsep=4pt}
\clubpenalty=10000 \widowpenalty=10000
\setcounter{tocdepth}{1}
\definecolor{focal}{HTML}{AF0032}
\definecolor{rose}{HTML}{FDF0F4}
\definecolor{border}{HTML}{D9DEE5}
\tikzset{box/.style={draw=border,rounded corners=2pt,line width=.55pt,
  text width=32mm,minimum height=11mm,align=center,inner sep=4pt,
  font=\small}, focus/.style={box,draw=focal,fill=rose},
  flow/.style={-{Latex[length=2mm]},line width=.65pt,draw=black!65}}

\title{LANCE : architecture du pipeline\\et méthodes de preuve}
\author{Codex, agent de diagnostic statique}
\date{Dernière révision : \revision\\[3pt]
  \small Exploration et protocoles proposés ; aucune expérimentation nouvelle}

\begin{document}
\maketitle
\begin{abstract}
Comment enrichir l'approche de LANCE avec une contribution vérifiable ? Une
lecture statique ciblée du code et des études existantes compare quatre
architectures : phases ordonnées, dossiers explicites de preuves et d'hypothèses,
comparaison différentielle de configurations, validation indépendante. Le code
contient déjà un registre de candidats, une attribution des appels d'outils et
une restitution déterministe. La difficulté est de garantir que chaque
conclusion conserve ses observations, ses contradictions et ses limites :
certains résultats d'outils sont filtrés avant leur journalisation. Les trois
extensions proposées peuvent d'abord être étudiées sur dossiers figés, sans
nouvelle collecte réseau. La copie de travail contient des conflits ; la lecture
ne démontre ni son exécutabilité, ni un gain de détection ou de fiabilité.
\end{abstract}
\tableofcontents

\section{Question et périmètre}
Une approche intéressante doit permettre de tester une proposition scientifique,
par exemple : une représentation explicite des preuves réduit-elle les
conclusions injustifiées et facilite-t-elle la revue humaine ? Le nombre
d'agents ou de phases est un choix d'architecture ; il ne constitue pas une
mesure de qualité de l'audit.

Ce rapport examine la circulation des observations et la formation des
conclusions dans LANCE, outil d'audit de réseaux de l'Internet des objets
(IoT). Il traite l'attribution, l'incertitude, les interruptions et l'isolation.
Il propose des études défensives sur données autorisées et déjà collectées ou
synthétiques. Aucun changement applicatif, appel au modèle LANCE, audit réseau,
exécution de laboratoire ou déploiement n'est réalisé.

Les recherches sur l'apport du modèle et la conservation du contexte ont déjà
proposé de comparer règles et modèle sur observations communes, avec des cas
sains et un coût suivi \cite{travaux}. Le présent diagnostic précise les points
d'appui actuels et distingue quatre variantes d'architecture. Il ne remplace
pas leurs protocoles détaillés.

\section{Méthode et état des sources}
\textbf{Lecture statique.} Les index \code{research/README.md} et
\code{docs/README.md}, les guides de phases et les études liées sont consultés,
puis les propriétaires du routage, de l'agrégation, du journal et du rapport.
Les constats ci-dessous décrivent les instructions visibles dans ces fichiers.
Les chemins et lignes permettent de revenir à chaque observation ; les numéros
peuvent changer lors d'une modification ultérieure \cite{code}.

\textbf{Version étudiée.} La base Git est
\code{a8dbeb9404f0cad960d61e42f02a1d70ca143a59}. La copie de travail contient
des modifications locales et des conflits non résolus, notamment dans
\code{src/agent/phases/verification/run.py}, \code{src/agent/exploit_evidence.py}
et plusieurs index. Le commit seul ne reproduit donc pas tous les textes lus.
Les branches en conflit ne sont ni choisies, ni fusionnées. Aucun module
applicatif n'est importé, aucun test applicatif exécuté et aucun comportement
du système déployé inféré de cette lecture.

\textbf{Sources externes.} Le résumé officiel du guide NIST SP~800-115 présente
la planification des examens, l'analyse des constats et la préparation de
mesures correctives ; seul ce résumé et sa notice sont lus ici \cite{nist}.
La note W3C \emph{PROV-Overview} décrit l'origine des données, les acteurs et
les transformations ; l'introduction et le guide de lecture sont consultés
\cite{prov}. Ce sont des appuis de vocabulaire. Ils ne démontrent pas que les
architectures proposées améliorent LANCE. La revue est ciblée, sans prétention
d'exhaustivité bibliographique.

\section{Constats : une base déjà structurée}
\subsection{Phases ordonnées et états distincts}
Le registre route six phases : graphe, reconnaissance, analyse, vérification,
évaluation des accès, rapport
(\src{src/agent/phases/registry.py}{11--26}). Le moteur les trie puis les
parcourt ; il contrôle arrêt, prérequis et conditions d'exécution
(\src{src/agent/pipeline.py}{663--753}). L'ordre des phases est donc visible
dans le code. Cela n'implique pas que tous les travaux internes soient
séquentiels : le scanner prévoit plusieurs tâches de travail selon le contexte.

Les états terminaux distinguent notamment arrêt, blocage, résultat partiel et
échec (\src{src/agent/results.py}{4--19}). Un état inconnu ne compte pas comme
succès. Ces états décrivent l'exécution ; la fin d'une phase ne prouve aucune
propriété de sécurité. De même, l'admission structurelle d'un candidat ne
certifie ni sa preuve ni sa compatibilité avec une vulnérabilité connue
(\src{src/agent/finding_schema.py}{1--45}).

\subsection{Provenance et conservation des informations}
Le registre de candidats conserve la valeur initiale, le fichier et le type de
source, l'indice, la décision et son motif. Il crée une copie de travail pour la
normalisation et garde les candidats mal formés avec leur rejet explicite
(\src{src/agent/phases/analysis/aggregation_loading.py}{20--51}). Une projection
brute distincte est annoncée dans l'agrégation
(\src{src/agent/phases/analysis/aggregation.py}{81--109}). C'est un point
d'appui pour expliquer pourquoi une proposition a été retenue ou écartée.

Le scanner marque ses appels \code{decision_source="rules"}
(\src{src/agent/phases/analysis/scan.py}{35--55}). Le journal commun enregistre
horodatage, durée, arguments, résultat, phase, agent, référence de preuve et
source de décision ; un échec d'écriture est explicite
(\src{src/agent/core/executor.py}{74--100}). Ce mécanisme permet de distinguer
une observation produite par une règle d'une proposition du modèle. Il ne
mesure pas, à lui seul, l'effet causal de chaque acteur sur le résultat final.

\textbf{Limite vérifiée.} Le chargeur d'outils peut filtrer des lignes et
tronquer la sortie avant de la retourner
(\src{src/agent/tools/tool_loader.py}{228--243}). L'exécuteur archive cette
valeur retournée (\src{src/agent/core/executor.py}{126--133}). Sur ce chemin,
le journal conserve donc une sortie déjà transformée, pas nécessairement la
capture native intégrale. Une référence stable ne permet pas de retrouver une
portion jamais conservée. L'effet réel sur les conclusions reste à mesurer.

\subsection{Frontières de restitution, d'arrêt et d'accès}
L'index du rapport distingue référence absente, ambiguë, incohérente avec
l'identité ou la destination. Il indique explicitement qu'il ne valide pas la
propriété de sécurité (\src{src/agent/phases/report/traceability.py}{1--69}).
Le rapport travaille par fiches, calcule sa synthèse et assemble les faits
déterministement. La rédaction d'une fiche est un appel sans outil ; sa cause
d'indisponibilité et les erreurs d'assemblage restent visibles
(\src{src/agent/phases/report/run.py}{57--118, 169--201}). Cette séparation
évite d'attribuer à une rédaction réussie la valeur d'une nouvelle preuve.

Le contrôle d'exécution vérifie arrêt et budgets avant les outils, puis refuse
de poursuivre un appel lorsqu'un arrêt est déjà visible
(\src{src/agent/core/executor.py}{30--49, 103--125}). Ces contrôles ne
démontrent pas l'annulation immédiate d'un processus déjà en cours.
Les chemins de livrables sont limités au dossier du run et plusieurs fichiers
internes sont masqués aux outils de l'agent
(\src{src/agent/artifacts.py}{6--67}). La documentation précise que ces
frontières ne sont pas un bac à sable système et que certains états, notamment
le graphe, restent partagés \cite{contrats}.

\begin{table}[htbp]
\centering\small
\caption{Acquis visibles et extensions proposées. Les acquis sont des constats
de lecture, pas des garanties vérifiées en exécution.}
\label{tab:faits}
\begin{tabularx}{\textwidth}{@{}p{29mm}YY@{}}
\toprule Sujet & Visible dans les sources lues & Extension à étudier \\
\midrule
Candidats & Valeur initiale, origine, décisions d'agrégation. & Relier également
les alternatives, contradictions et inconnues. \\
Preuves & Références et sorties retournées dans un journal. & Conserver la
capture avant projection, ou déclarer précisément son absence. \\
Attribution & Source de décision et agent enregistrés. & Comparer des interprètes
sur les mêmes dossiers ; distinguer collecte et interprétation. \\
Restitution & Faits assemblés par le code, fiches sans outils. & Rendre chaque
conclusion contestable par un dossier de preuves lisible. \\
Exécution & États, limites, dossiers de livrables et refus explicites. & Vérifier
l'arrêt en cours d'appel et l'isolation de tous les états par tests dédiés. \\
\bottomrule
\end{tabularx}
\end{table}

\section{Quatre architectures à comparer}
\subsection{Phases ordonnées : référence de comparaison}
La première variante conserve le chemin actuel et les contrats de preuve. Son
intérêt est de fournir une référence compréhensible : observations, candidats,
vérification, restitution. Sa limite est qu'un livrable compact peut omettre
une alternative utile à une étape ultérieure. Le diagnostic ne démontre pas
que cette perte se produit sur tous les cas. Avant comparaison, il faut figer
une version exécutable résolue et documenter quelles informations chaque phase
reçoit effectivement. Ajouter une phase sans préciser ce qu'elle mesure ne
constituerait pas une nouvelle contribution scientifique.

\subsection{Dossiers de preuves et d'hypothèses explicites}
\textbf{Proposition.} Chaque propriété examinée reçoit un dossier : énoncé
précis, objet concerné, observations pertinentes, origine et date, éléments en
faveur et en contradiction, inconnues, règle de conclusion et historique des
révisions. Une observation décrit ce qui a été recueilli ; une hypothèse décrit
ce qu'on en déduit. La conclusion peut rester indéterminée. Le modèle explique
ces dossiers, tandis que les règles retenues restent responsables de
l'acceptation des preuves.

Cette variante prolonge le registre existant. Le vocabulaire de provenance
de W3C PROV suggère de représenter séparément données, transformations et
acteurs \cite{prov}. Cette adaptation à LANCE est une proposition de ce rapport,
sans revendication de conformité PROV. Il faut notamment conserver les
observations avant projection, indiquer leur caractère complet ou partiel et
garder une vue compacte qui renvoie à ce qui existe réellement.

La figure~\ref{fig:dossiers} explicite une distinction centrale : une
validation peut contester une conclusion sans modifier ses observations.

\begin{figure}[htbp]
\centering
\begin{tikzpicture}
\node[box] (obs) at (0,0) {Observations\newline figées};
\node[focus] (dossier) at (5,0) {Hypothèses et\newline preuves liées};
\node[box] (report) at (10,0) {Conclusions\newline et limites};
\node[box] (review) at (5,-2.0) {Validation\newline indépendante};
\draw[flow] (obs) -- (dossier);
\draw[flow] (dossier) -- (report);
\draw[flow] (obs.south) |- (review.west);
\draw[flow] (dossier.south) -- (review.north);
\draw[flow] (review.east) -| (report.south);
\end{tikzpicture}
\caption{Architecture proposée de restitution sur données figées. Le validateur
lit les observations et les hypothèses, puis fournit une appréciation distincte ;
il ne réécrit pas les observations. Aucune collecte réseau n'est représentée.}
\label{fig:dossiers}
\end{figure}

\subsection{Audit différentiel : configurations fragiles et corrigées}
\textbf{Proposition.} Comparer deux dossiers issus de configurations contrôlées
dont une propriété défensive varie : par exemple un inventaire ou une
configuration de politique d'accès avant et après correction. Dans la première
étude, utiliser des fichiers de configuration fournis et des observations
synthétiques clairement déclarées ; aucune reproduction de faille n'est
nécessaire. Les dossiers sont codés sans indiquer au classificateur lequel est
le cas attendu sain. L'annotation de référence est tenue à part.

La question porte sur la cohérence : la conclusion relative à la propriété
modifiée change-t-elle, tandis que les conclusions relatives aux propriétés
inchangées restent stables ? Un résultat différent ne démontre une relation
causale que si les autres changements pertinents sont maîtrisés. Un cas dit
« corrigé » peut présenter d'autres risques ; une différence de bannière ou
de nom ne doit pas suffire à changer le verdict. Cette variante mesure une
sensibilité à une correction connue et les fausses alertes, au-delà du nombre
brut de constats produits.

\subsection{Validation indépendante des conclusions}
\textbf{Proposition.} Une seconde lecture reçoit les observations originales et
les règles d'acceptation gelées. Elle recherche les conclusions sans preuve,
les mauvaises attributions et les contradictions non résolues. Son appréciation
est publiée séparément ; un désaccord ouvre un point à revoir, sans supprimer
silencieusement le candidat ni inventer une confirmation. Une revue humaine
peut fournir une annotation indépendante pour le corpus étudié.

L'indépendance exige de préciser ce qui est partagé : observations, modules d'interprétation,
taxonomie, règles et éventuellement famille de modèle. Deux agents utilisant
le même parseur et la même règle erronée peuvent s'accorder. Un vote entre
modèles ne remplace donc pas une référence défensive documentée. La valeur
recherchée est la détection d'erreurs de conclusion ; l'accord seul n'est pas
la métrique de réussite.

\begin{table}[htbp]
\centering\small
\caption{Ce que chaque comparaison pourrait établir, sous conditions. Aucun
gain n'est observé dans cette étude.}
\label{tab:architectures}
\begin{tabularx}{\textwidth}{@{}p{31mm}YY@{}}
\toprule Variante & Question mesurable & Condition principale \\
\midrule
Phases ordonnées & Quel est le comportement de référence ? & Version résolue,
entrées et sorties effectivement figées. \\
Dossiers explicites & Les conclusions restent-elles reliées à leurs preuves ? &
Mêmes observations et règles ; seule la représentation varie. \\
Différentielle & La réponse suit-elle une correction défensive connue ? &
Variation contrôlée, annotation séparée et cas non concernés. \\
Validation indépendante & Des erreurs de conclusion sont-elles détectées ? &
Référence de contrôle indépendante et dépendances partagées explicites. \\
\bottomrule
\end{tabularx}
\end{table}

\section{Hypothèses et protocole proposé}
L'étude minimale peut être réalisée hors ligne. Elle porte sur la qualité des
conclusions et de leur justification ; elle ne mesure pas l'exploration réseau.
Trois hypothèses sont réfutables : les dossiers explicites réduisent les
affirmations sans preuve ; les paires différentielles révèlent des changements
de conclusion incohérents ; la validation séparée repère plus d'erreurs que la
seule vérification de références. Aucun de ces effets n'est acquis.

Le protocole proposé comporte quatre étapes :
\begin{enumerate}
\item \textbf{Fixer les observations.} Préparer un corpus de dossiers autorisés
et existants, ou de dossiers synthétiques déclarés, incluant cas corrects,
preuves manquantes, contradictions et mauvaises attributions. Consigner la
provenance, les portions absentes et la version de chaque fichier.
\item \textbf{Fixer la référence.} Faire annoter propriétés et exigences de
preuve avant comparaison ; résoudre les désaccords humains et conserver les
cas indéterminés. Une simple sortie attendue du code existant ne fournit pas
une validation indépendante de sa sémantique.
\item \textbf{Varier une dimension.} Comparer représentation actuelle et dossiers
explicites avec le même interprète. Ajouter ensuite la lecture indépendante,
puis les paires défensives. Inclure un interprète à règles et, si le coût et
l'accès le permettent, un modèle de langage sans outil. Garder constantes les
autres entrées et les exigences de preuve.
\item \textbf{Publier erreurs et coûts.} Compter les conclusions injustifiées,
les liens de preuve exacts, les contradictions oubliées, les fausses alertes
et les abstentions. Mesurer aussi temps de revue, coût modèle et issues
incomplètes. Publier les dénominateurs et valeurs manquantes ; les instances
partageant un même mécanisme ne sont pas des observations indépendantes.
\end{enumerate}

Une amélioration de liens de preuve peut accompagner une hausse des abstentions.
Il faut donc définir avant comparaison ce qui constitue une amélioration utile,
et publier cette contrepartie. Des données synthétiques réussies justifieraient
au plus la poursuite sur des dossiers réels autorisés ; elles ne démontreraient
pas la qualité d'un audit en production.

\section{Limites et conclusion}
La lecture ne fournit aucun taux d'erreur, coût expérimental ou temps humain
mesuré. Les constats de code sont limités aux chemins inspectés ; les modules
en conflit empêchent de présenter cette copie comme une référence exécutable.
Les études antérieures décrivent leurs propres dates et expériences ; leurs
résultats ne sont pas rejoués ici. Les architectures proposées n'ont été ni
implémentées, ni comparées.

Une garantie complète d'isolation demanderait de vérifier état du graphe,
processus, fichiers, entrées publiques et journaux. Une interruption fiable
demanderait de distinguer refus d'un prochain appel, arrêt d'un appel en cours
et finalisation avec les données déjà disponibles. Ces objectifs restent des
conditions de validation logicielle ; une réussite aux tests ne prouverait pas
une amélioration scientifique de l'audit.

\textbf{Conclusion de conception.} LANCE dispose déjà de frontières pertinentes
entre candidats, observations, exécution et restitution. La piste la plus
directement étayée par ce diagnostic est de rendre les dossiers de preuve
explicites et récupérables, puis d'évaluer la fidélité des conclusions sur
observations communes. La comparaison différentielle et la validation
indépendante apportent deux contrôles complémentaires. Leur intérêt reste une
hypothèse mesurable. Le prochain travail utile est de figer les règles, le corpus
et les critères de jugement avant toute affirmation de gain.

\begin{thebibliography}{9}
\addcontentsline{toc}{section}{Références}
\bibitem{travaux} Recherches LANCE, \emph{Agent / automatisation} et
\emph{Conservation du contexte}, rapports révisés le 1er octobre 2026 à 10:23,
Europe/Brussels. Sources et analyses ciblées consultées le 7 octobre 2026 ;
expériences antérieures non rejouées.
\href{../2026-09-27-agent-vs-automation/README.md}{Dossier agent / automatisation} ;
\href{../2026-09-29-context-scalability/README.md}{dossier contexte}.

\bibitem{code} Dépôt LANCE, copie de travail inspectée le 7 octobre 2026,
base \code{a8dbeb9404f0cad960d61e42f02a1d70ca143a59}, modifications locales et
conflits non résolus. Les liens de fichiers et lignes dans les constats sont
la référence de lecture ; aucun résultat d'exécution nouvelle.
\href{../../src/agent/phases/README.md}{Guide des phases},
\href{../../docs/README.md}{index documentaire} et
\href{../README.md}{conventions de recherche}.

\bibitem{contrats} Documentation LANCE,
\href{../../docs/architecture/run-artifacts.md}{\emph{Livrables : un dossier
explicite par run}} et
\href{../../docs/architecture/report-writing.md}{\emph{Rapport rédigé par
sections}}. Pages de la copie de travail consultées le 7 octobre 2026, pas une
attestation de validation actuelle ni de déploiement.

\bibitem{nist} K. Scarfone, M. Souppaya, A. Cody et A. Orebaugh,
\emph{Technical Guide to Information Security Testing and Assessment},
NIST SP 800-115, septembre 2008.
\url{https://csrc.nist.gov/pubs/sp/800/115/final}.
Notice et résumé officiel lus le 7 octobre 2026 ; guide complet non dépouillé.

\bibitem{prov} P. Groth et L. Moreau (éd.), \emph{PROV-Overview: An Overview
of the PROV Family of Documents}, W3C Working Group Note, 30 avril 2013.
\url{https://www.w3.org/TR/2013/NOTE-prov-overview-20130430/}.
Introduction et guide de lecture consultés le 7 octobre 2026 ; aucune
implémentation ni validation de conformité PROV réalisée.
\end{thebibliography}
\end{document}
